% CAS (Heimdall) / SiteControl -- requirements for OPC UA clients of the PLC.
% Build: latexmk -pdf -interaction=nonstopmode "OPC UA client requirements.tex"
\documentclass[11pt,a4paper]{article}
\usepackage{amsmath}
\makeatletter
\def\input@path{{Functional description/}}   % casfd.sty lives with the FDs
\makeatother
\usepackage{casfd}
\setlength{\emergencystretch}{2em}

% Requirement table: ID, requirement, rationale / source
\newenvironment{reqtable}{%
  \begin{small}
  \setlength{\tabcolsep}{4pt}%
  \begin{longtable}{@{}P{14mm}P{86mm}P{60mm}@{}}
  \toprule
  ID & Requirement & Why / source\\
  \midrule
  \endfirsthead
  \toprule
  ID & Requirement & Why / source\\
  \midrule
  \endhead
  \bottomrule
  \endfoot
}{%
  \end{longtable}
  \end{small}
}
\newcommand{\TBD}{\textcolor{red}{TBD}}

\title{CAS (Heimdall) -- requirements for OPC UA clients of the PLC\\[1mm]\Large HMI, SiteControl and tools}
\author{ScaleAQ Automation \quad Draft, rev.~A0}
\date{\today}

\begin{document}
\maketitle

\begin{abstract}
\noindent Every program that talks OPC UA to the CAS PLC shares one small controller with the control
program. This document states what such a client must do, and must not do, so that it gets reliable
data and does not load or disturb the PLC. It applies to the CAS web HMI (\var{heimdall-cas-hmi}),
to SiteControl services that read or command the PLC, and to engineering tools. The requirements come
from the CAS commissioning on the simulator (November 2026), where a client reading every value every
second turned out to use about a third of the PLC's CPU. Items not decided are marked \open{n} and
listed in Section~\ref{sec:open}.
\end{abstract}

\tableofcontents

% =============================================================================================
\section{Scope and the server}
\label{sec:scope}

\subsection{Clients in scope}
\begin{itemize}[nosep]
  \item \textbf{CAS web HMI} (\var{scaleaq/heimdall-cas-hmi}, Python \var{asyncua}): reads about 600 values,
        writes operator commands and alarm acknowledges.
  \item \textbf{SiteControl services} that integrate the CAS PLC (edge drivers on
        \var{sitecontrol-opcua-lib}, the OPC Foundation .NET stack): read plant values and alarms,
        possibly send commands.
  \item \textbf{Engineering tools}: UaExpert, the CODESYS IDE (online monitoring, trace), used during
        commissioning and service only.
\end{itemize}
Not in scope: the PLC as an OPC UA \emph{client} (its data sources to the Framo pump controllers), and
Modbus clients.

\subsection{The server}
\begin{itemize}[nosep]
  \item WAGO PFC200 G2 (750-8212), \textbf{one CPU core}, 512\,MB RAM, CODESYS runtime with its OPC UA
        server on port 4840. The same core runs the control task (20\,ms), the I/O, the Modbus masters and
        the PLC's own OPC UA client.
  \item The address space is the application's \emph{Symbol Configuration}: \var{GVLMain}
        and the variables ticked there. Node ids are CODESYS string ids,
        \var{ns=4;s=|var|<device>.Application.<GVL>.<path>}; device name and namespace index differ
        between projects.
  \item The real PLC requires a user login; anonymous access is off.
  \item A download, Reset Origin or reboot restarts the address space. An online change can add
        variables without ending existing sessions.
\end{itemize}

\subsection{What the measurements showed}
\label{sec:measure}
Simulator PLC 10.1.2.41, 2026-11-03, \var{htop} on the controller. The HMI read 455 values and 28
quality flags every second and 136 data log values every 2\,s, each by string node id, in requests of 50.

\begin{center}
\begin{small}
\begin{tabular}{@{}lccc@{}}
\toprule
State & Runtime main thread & \var{OPCUAServer} thread & Total CPU\\
\midrule
No application loaded & $\approx$15\,\% & $\approx$5\,\% & 26\,\%\\
Application running, HMI blocked & 40\,\% & $\approx$0 & 33\,\%\\
Application running, HMI reading & 55--65\,\% & 21--25\,\% & 66--68\,\%\\
Application stopped, HMI reading & 53\,\% & 25\,\% & 57\,\%\\
\bottomrule
\end{tabular}
\end{small}
\end{center}

The cost followed the client, not the program: it stayed with the application stopped, and disappeared
when the client was blocked. The server resolves a string node id on every read. A subscription lets it
resolve each variable once (\var{heimdall-cas-hmi} PR~\#25, to be measured).

% =============================================================================================
\section{Requirements}
\label{sec:req}
``Shall'' is a requirement, ``should'' a recommendation.

\subsection{Connection and session}
\begin{reqtable}
CON-01 & The client shall log in with its own user account on the PLC. Credentials shall come from the
deployment's secrets (environment, Kubernetes secret), never from source code or images. &
Real PLC has no anonymous access. One account per client makes a misbehaving client identifiable and
lets it be locked out alone.\\
CON-02 & The client shall use one session per PLC and give it a session name that identifies the client
and its site. & Several sessions per client multiply the server load; the name shows up in the
runtime log.\\
CON-03 & After a lost connection the client shall reconnect by itself, with a back-off that starts at
about 1\,s and grows to at most 15\,s between attempts. & Fast reconnect after a download; no connection
storm while the PLC reboots, and back within 15\,s once it is up.\\
CON-04 & On every new session the client shall look up its nodes again, and shall not reuse node ids
from an earlier session. & Download, Reset Origin and reboot rebuild the address space.\\
CON-05 & A node that is not found shall be reported once (log, and ``not available'' where it is shown),
shall not stop the client, and shall be looked up again at least every 10 minutes and at every reconnect.
& Symbols added by an online change only appeared in the HMI after its pod was restarted.\\
CON-06 & The client shall not hard-code the namespace index or the device name. It should find the GVL
root (\var{GVLMain}) once per session by browsing and derive the other paths from it. & Device and
application names differ between the simulator PLC, the site PLC and the soft PLC.\\
CON-07 & Security policy and certificates \open{1}. & Not decided for site use.\\
\end{reqtable}

\subsection{Reading data}
\begin{reqtable}
RD-01 & Cyclic data shall be read through a \textbf{subscription} (monitored items), not by reading every
value every cycle. & Section~\ref{sec:measure}: polling 600 values per second used about a third of the CPU.\\
RD-02 & The sampling and publishing interval shall be at least 1\,s. Faster sampling shall be limited to
a few values, only while a user looks at them, and at most 5\,Hz. & The HMI's wave trend samples two
level transmitters at 5\,Hz only while a page shows it.\\
RD-03 & Monitored items shall use queue size 1 (latest value only). & The client needs the current state,
not a history; queues cost server memory.\\
RD-04 & A client that must read (no subscription available, or a value the server refuses to monitor)
shall read at most 50 values per request and should not read the same value more often than every 2\,s.
& The PFC200 refuses large requests (\var{BadTooManyOperations} at 200).\\
RD-05 & Large browses (for example all \var{AlarmAck} bits, about 1100 nodes, 35\,s) shall run at most
once per session, in the background, and never periodically. & Browsing walks the symbol tree in the
runtime.\\
RD-06 & The client shall read only what it uses. Each client shall keep its list of paths in one place
in its code, with the PLC data type. & The list is the contract with the PLC (Section~\ref{sec:contract}).\\
RD-07 & Total load: with all permanent clients connected the PLC shall keep \TBD{} CPU in reserve
\open{2}. & Measured baseline 33\,\% without clients.\\
\end{reqtable}

\subsection{Link supervision and data quality}
\begin{reqtable}
LQ-01 & The client shall notice a lost connection within 3 publishing intervals (keep-alive, or a cheap
read of \var{Server_ServerStatus_CurrentTime}, a numeric node). & A subscription alone sends nothing
when nothing changes.\\
LQ-02 & While the connection is lost, every value shall be shown as unknown, never as the last value
received. & An operator must not act on frozen data.\\
LQ-03 & The client shall respect the PLC's own quality information: values whose quality flag says bad
(\var{Status.BadQuality}), and values from a device whose link the PLC reports down or stale
(\var{CommOK} FALSE, \var{DataStale} TRUE), shall be shown as invalid. & The PLC knows more about its
field devices than the client.\\
LQ-04 & The client shall show what the PLC computes (scaling, averages, states) and shall not compute
plant values itself. & One source of truth; the FDs describe the PLC.\\
LQ-05 & Timestamps: the client shall use its own clock for when it received a value until the PLC
clock is synchronised \open{3}. & The simulator PLC clock was months off.\\
\end{reqtable}

\subsection{Writing}
\begin{reqtable}
WR-01 & The client shall write only variables on its own whitelist: \var{Command} members, \var{AlarmAck}
bits and named setpoints. It shall never write \var{Status}, \var{Input}, \var{Output},
\var{InternalCommand} or \var{Alarm}. & The PLC owns its state.\\
WR-02 & Commands are pulses: write TRUE, hold at least 1\,s or until the PLC clears it, then FALSE.
The client shall not rely on the PLC clearing a bit it did not document as cleared. & AlarmManagement
acknowledges on the rising edge; some commands are latched and cleared by the PLC.\\
WR-03 & On connect and on a clean session end, the client shall write FALSE to every command bit it may
have left TRUE. & A command bit left TRUE after a lost link would act again on the next edge.\\
WR-04 & Setpoints shall be bounds-checked by the client before writing, with the same limits as the PLC.
& The PLC limits them too; the client gives the operator a clear message instead.\\
WR-05 & Retained configuration (\var{InletConfig}, \var{OxyConfig}, IP addresses) shall only be written by
explicitly named, whitelisted functions, never in bulk. & Retained values survive a restart; a wrong bulk
write survives too.\\
WR-06 & Every client shall have a read-only mode that writes nothing. & Diagnosis on a running plant.\\
WR-07 & Only one client shall command a given function at a time \open{4}. & HMI and SiteControl must not
fight over the same command.\\
\end{reqtable}

\subsection{The contract with the PLC}
\label{sec:contract}
\begin{reqtable}
CT-01 & Everything a client reads or writes shall be in the Symbol Configuration. New fields of a
structure shall be ticked there explicitly (where the structure is not exported whole). & New fields
were in the PLC but not on the server until ticked.\\
CT-02 & A PLC change that renames, removes or retypes a published variable shall update the clients'
path lists in the same change, or be agreed with the client owners first. & The client finds out only at
runtime.\\
CT-03 & The PLC should publish a program version (\var{PlcInfo}) that clients log at connect. & To tell
which program a client talked to.\\
\end{reqtable}

\subsection{Engineering tools}
\begin{reqtable}
EN-01 & UaExpert, the CODESYS IDE online view and traces shall be connected only while needed, and
closed after commissioning and service. & They add load (subscriptions, monitoring) on the same core.\\
EN-02 & CPU checks on the controller shall be done with the IDE logged out and UaExpert closed,
reading \var{htop} with thread names (\var{OPCUAServer}, \var{codesys3}). & Otherwise the tools
measure themselves.\\
\end{reqtable}

% =============================================================================================
\section{Verification}
\label{sec:verify}
A client is accepted when, on the simulator PLC with the application running:
\begin{enumerate}[nosep]
  \item total CPU with the client connected is within \TBD{} of the value with the client blocked
        (\var{iptables -I INPUT -p tcp --dport 4840 -s <client IP> -j REJECT}, removed afterwards) \open{2};
  \item it notices a lost link within 3 publishing intervals and shows its values as unknown;
  \item it reconnects by itself after a download and after a reboot, and finds its nodes again;
  \item it finds a variable added by an online change without a restart (CON-05);
  \item it writes nothing in read-only mode.
\end{enumerate}

% =============================================================================================
\section{Open issues}
\label{sec:open}
\begin{small}
\begin{longtable}{@{}P{10mm}P{86mm}P{58mm}@{}}
\toprule
No. & Issue & To decide\\
\midrule
\endhead
\bottomrule
\endfoot
OI1 & Security policy for site use: None, or Basic256Sha256 with certificates; who manages them. &
Automation / IT.\\
OI2 & CPU budget: how much the PLC shall keep in reserve with all clients connected, and how much one
client may use. & After the subscription has been measured (PR~\#25).\\
OI3 & NTP on the PLC; until then source timestamps are not trusted. & Commissioning.\\
OI4 & Which client commands what when both the HMI and SiteControl are connected. & Project.\\
OI5 & Does SiteControl read the CAS PLC directly over OPC UA, or through the HMI or another gateway? &
SiteControl team.\\
\end{longtable}
\end{small}

\end{document}
